\chapter{相关概念与技术基础}

本论文所提出的多签名在网检测（第三章）与多星协同推回缓解（第四章）方法，其设计与实现深度依赖于四个领域的理论支撑：大规模 LEO 星座的网络拓扑与约束模型、链路洪泛攻击的构造机理与对抗博弈、P4 可编程数据面的架构原理与计算边界、以及分布式 Gossip 协议的传播理论与开销控制。本章将逐一展开论述，为后续章节建立可直接引用的理论基础。

\section{大规模低轨卫星网络体系架构与特性}

\subsection{巨型星座的网络拓扑与星间链路模型}

现代 LEO 巨型星座的网络拓扑可从两个层面进行形式化描述：物理构型层面与网络图论层面。

在物理构型层面，巨型星座的卫星分布遵循 Walker 星座模型。以 Walker Delta 构型为例（如 Starlink 第一阶段主体），其参数化表示为 $W(\delta: T/P/F)$，其中 $\delta$ 为轨道倾角，$T$ 为卫星总数，$P$ 为轨道平面数，$F$ 为相邻轨道间的相位因子。对于 Starlink Shell 1（$\delta=53°, T=1584, P=72, F=1$），每个轨道平面承载 $T/P = 22$ 颗卫星，轨道高度 $h = 550$\,km。Walker Star 构型（如 Telesat Lightspeed，$\delta=99.5°$，近极地轨道）则在极地地区存在轨道交叉区，形成拓扑上的特殊汇聚点。两种构型的共同特征是：在同一轨道平面内，相邻卫星之间通过同轨星间链路（Intra-plane ISL）相连，链路长度相对稳定；不同轨道平面之间，通过异轨星间链路（Inter-plane ISL）相连，其长度随卫星纬度位置的变化而周期性变化~\cite{laser_isl}。在高纬度区域（特别是极圈附近），异轨 ISL 由于相邻轨道平面间距过大或卫星相对速度过高，可能周期性中断~\cite{walker1984satellite, kodheli2021satellite}。

在网络图论层面，LEO 星座在任意给定时刻 $t$ 的网络拓扑可建模为时变有向图 $G(t) = (V(t), E(t))$，其中 $V(t)$ 为当前在轨运行的卫星节点集合，$E(t)$ 为当前活跃的 ISL 集合。对于典型的四端口 ISL 设计（前后两条同轨 ISL、左右两条异轨 ISL），除极地交叉区外，每个卫星节点的网络度数约为 4。这种规则的格状（Grid/Torus）拓扑结构意味着：（1）任意两颗卫星之间的最短路径跳数可由其在网格中的曼哈顿距离（Manhattan Distance）近似估计；（2）当某条核心 ISL 发生拥塞时，替代路径的跳数增量相对可控，但同时也意味着拥塞可能沿网格的行或列方向传播，引发连锁效应~\cite{handley2019delay, icarus}。

每条 ISL $e \in E(t)$ 具有有限的链路容量 $C_e$（通常为 10$\sim$100\,Gbps 量级）。定义链路 $e$ 在时刻 $t$ 的利用率为 $\rho_e(t) = \lambda_e(t) / C_e$，其中 $\lambda_e(t)$ 为 $e$ 上的瞬时聚合流量速率。当 $\rho_e(t) > 1$ 时，链路进入拥塞状态，数据包在输出缓冲队列中排队等待或被丢弃。LFA 攻击的目标正是通过人为制造 $\rho_{e^*}(t) > 1$（$e^* \in E_{\text{target}}$）来使目标链路 $e^*$ 瘫痪~\cite{crossfire}。这一量化模型为第三章的检测阈值设定和第四章的推回触发条件提供了形式化基础。


% =========================
% 扩展块 A：拓扑规则性与安全脆弱面的联系
% 建议插入位置：第 15 行后
% =========================
从网络安全视角看，Walker 型巨型星座并不是一张“各链路风险均匀分布”的理想网络。尽管同轨与异轨 ISL 在物理构型上呈现高度规则性，但在业务需求分布、信关站部署、低时延路由偏好以及极区链路可用性波动的共同作用下，不同链路被端到端路径选中的概率并不相同。若将某一快照时刻的活跃端到端路径集合记为 $\mathcal{P}(t)$，则可定义链路 $e$ 的路径承载比为
\begin{equation}
    \eta_e(t)=\frac{1}{|\mathcal{P}(t)|}\sum_{p\in\mathcal{P}(t)} \mathbf{1}(e\in p),
\end{equation}
其中 $\mathbf{1}(\cdot)$ 为指示函数。对长期处于高 $\eta_e(t)$ 区域的链路而言，即便其物理容量并未低于其他链路，也会因路径汇聚效应而更容易成为业务热点和攻击目标。换言之，LEO 星座中的“关键链路”并不一定是物理层面最稀缺的链路，而更可能是路由层面被大量最短路或低时延路偏好的链路~\cite{handley2019delay, icarus, giuliari2020internet_backbones}。

这一观察对本文后续研究有两层直接含义。第一，第三章的单星检测不应仅把链路拥塞看作孤立现象，而应将其理解为“路径汇聚结构”在本地节点上的投影；第二，第四章的协同区域（ROI）也不宜脱离路径热区独立选取，而应围绕高汇聚链路的上下游卫星群展开。也正因为如此，第二章不仅需要解释 LEO 星座“长什么样”，更需要解释“为何某些链路会反复成为风险焦点”，这正是从网络拓扑描述走向安全建模描述的关键一步。

% 建议新增图 2-1（待绘制后插入）
% 图题建议：Walker 巨型星座的物理构型与逻辑网格抽象
% 画法建议：左侧给出 Walker Delta / Walker Star 两种典型构型，标出同轨与异轨 ISL；
%            右侧将其抽象为 4-regular mesh，并用加粗红色链路标出高路径汇聚链路。
% \begin{figure}[t]
%   \centering
%   \includegraphics[width=0.92\linewidth]{figs/ch2/fig2_constellation_arch.pdf}
%   \caption{Walker 巨型星座的物理构型与逻辑网格抽象。}
%   \label{fig:ch2_constellation_arch}
% \end{figure}

\subsection{卫星网络的高动态性与快照路由模型}

LEO 卫星以约 7.5\,km/s 的速度运行（对于 $h=550$\,km 的圆轨道，轨道周期约 $T_{\text{orbit}} = 2\pi\sqrt{(R_E+h)^3/\mu} \approx 95.6$\,min，其中 $R_E = 6371$\,km 为地球半径，$\mu = 3.986 \times 10^{14}\,\text{m}^3/\text{s}^2$ 为地心引力常数）。这一高速运动带来两个层面的动态性~\cite{kodheli2021satellite}。

第一层面为星地覆盖的动态性：每颗卫星的地面覆盖区（Footprint）持续移动，地面终端需要每隔数分钟执行一次卫星切换（Satellite Handover）。第二层面为星间拓扑的动态性：异轨 ISL 的物理可用性和链路质量随卫星间相对位置的变化而波动；在极地区域，异轨 ISL 的周期性中断导致网络在该区域的连通性呈现间歇性特征。

为应对上述动态性，学术界广泛采用离散时间快照（Discrete-Time Snapshot）模型。其核心思想是利用 LEO 轨道运动的完全可预测性（由开普勒轨道力学确定），将连续时间轴划分为一系列等长的离散时段 $\{[\tau_k, \tau_{k+1}) | k=0,1,2,\ldots\}$，时段长度 $\Delta\tau = \tau_{k+1} - \tau_k$（典型值为 10$\sim$30\,s，如 Starlink 约15\,s）。在每个时段 $[\tau_k, \tau_{k+1})$ 内，网络拓扑 $G(\tau_k)$ 被视为静态不变，卫星在时段开始时加载由地面网络管理系统预先计算并上注的转发信息表（Forwarding Information Base, FIB）。时段结束时，所有路由状态同步更新至下一时段的 FIB~\cite{handley2019delay, hypatia}。

% [INS-ICARUS-04]
ICARUS 的攻击建模进一步说明了该快照假设的安全含义：在链路最短路变化以秒级发生、而单次端到端转发时延远小于快照长度的条件下，攻击者可将连续时变系统近似为一系列静态快照，并提前计算多时段攻击流分配；因此“拓扑动态性”并不必然抬升攻击门槛，关键在于防御端是否能在同等时间尺度完成检测与响应闭环~\cite{icarus}。

快照模型对本文的防御方案设计具有两方面的直接意义。其一，在第三章中，在网检测算法的状态累积周期（检测窗口 $W$）可以与快照时段 $\Delta\tau$ 对齐——在每次拓扑切换时，对检测状态表进行优雅的清洗（Graceful Reset），既避免了跨拓扑期间的状态污染，又保证了窗口内统计的一致性。其二，在第四章中，由于同一快照时段内的路由路径是确定的，推回令牌（Pushback Token）可以沿着已知的反向路径精确回传，无需实时的路径发现机制。


% =========================
% 扩展块 B：快照模型的安全意义与多时间尺度分离
% 建议插入位置：第 28 行后
% =========================
\paragraph{快照模型的安全意义}
快照路由模型之所以能够成为 LEO 网络研究中的主流抽象，并不仅仅因为它便于仿真，更重要的是它准确刻画了该类网络中“多时间尺度分离”的系统事实：单个数据包的端到端转发时延通常处于毫秒级，检测窗口和局部协同过程多位于数百毫秒到数秒量级，而由轨道运动驱动的拓扑显著变化则常发生在秒级到十秒级。因而在多数工程场景下可近似满足
\begin{equation}
    T_{\mathrm{pkt}} \ll W \lesssim \Delta\tau \ll T_{\mathrm{orbit}},
\end{equation}
其中 $T_{\mathrm{pkt}}$ 表示包级转发时间尺度，$W$ 表示检测或聚合窗口，$\Delta\tau$ 表示快照长度，$T_{\mathrm{orbit}}$ 表示轨道周期。该不等式意味着：在单个快照内部，防御系统通常可以把转发路径视为近似静态，从而在局部完成状态累积、异常判定与缓解动作，而不必在每个数据包层面重新面对拓扑不确定性~\cite{hypatia, handley2019delay, icarus}。

然而，快照模型也带来两个常被忽略的安全问题。其一，若检测窗口 $W$ 明显大于 $\Delta\tau$，则窗口中的流量统计可能跨越两套不同转发表，导致“同一 key 的统计意义”在窗口内部发生漂移，破坏阈值判决的稳定性。其二，若 FIB 切换并非所有卫星同时完成，而是存在短暂的不一致阶段，则攻击者可能利用切换边界制造瞬时的路径重分布，使局部观测进一步碎片化。正因如此，本文后续设计中强调将检测窗口、位图双缓冲翻转、协同证据寿命和推回软状态寿命尽量与 epoch/快照边界对齐，以减少“状态属于旧拓扑，而策略作用于新拓扑”的错位现象。

从论文写作角度看，这一段内容值得在第二章中明确展开，因为它能够帮助读者理解：为什么第三章与第四章大量采用 epoch、窗口、快照对齐等设计；为什么“高动态”并不意味着“完全无法建模”；以及为什么攻击者与防御者都可能在同一快照框架下展开博弈。换句话说，快照模型不是简单的背景假设，而是本文算法设计与安全分析的共同时间基准。

% 建议新增图 2-2（待绘制后插入）
% 图题建议：LEO 快照路由、检测窗口与拓扑切换的时间尺度关系
% 画法建议：横轴为时间，画出连续的 snapshot k / k+1 / k+2，叠加 FIB 版本、检测窗口 W、
%            协同传播窗口和推回软状态寿命，突出“对齐”与“错位”两种情况。
% \begin{figure}[t]
%   \centering
%   \includegraphics[width=0.96\linewidth]{figs/ch2/fig2_snapshot_timeline.pdf}
%   \caption{快照路由、检测窗口与拓扑切换的时间尺度关系。}
%   \label{fig:ch2_snapshot_timeline}
% \end{figure}

\subsection{面向星上处理的软硬件资源约束}

将安全检测与缓解逻辑部署于卫星平台，必须正视航天工程的 SWaP（Size, Weight, and Power）约束对计算与存储资源施加的刚性上限~\cite{satshield, kodheli2021satellite}。

在功耗维度，LEO 卫星的电力供应完全依赖太阳能电池阵列。以典型的中型通信卫星为例，其太阳能帆板的发电功率约为 1$\sim$5\,kW，而这一功率预算需在通信载荷（星载天线、射频前端）、计算单元、姿态控制系统（ADCS）和热控系统之间分配。分配给网络交换与处理单元的功耗预算通常不超过数十瓦。作为对比，地面数据中心的高性能可编程交换芯片（如 Intel Tofino 2, 12.8\,Tbps）的典型功耗约为 200$\sim$300\,W，仅此一项即可能超出整颗小型卫星的总功率预算~\cite{p4_survey, kodheli2021satellite}。

在存储维度，地面 Tofino 2 芯片提供约 80\,MB 的片上 SRAM 和 40\,Mb 的 TCAM。而宇航级 FPGA（如 Xilinx Virtex-5QV 系列抗辐射器件）的片上 BRAM 容量通常仅为数 MB 量级。如果以 Per-flow 粒度维护流状态表（假设每条流需要 $\sim$64\,Bytes 的状态信息，活跃流数量为 $N_f$），则精确状态维护所需的内存为 $M_{\text{exact}} = 64 N_f$\,Bytes。当 $N_f = 10^6$（中等规模流量）时，$M_{\text{exact}} = 64$\,MB，这已远超宇航级芯片的片上SRAM容量~\cite{p4_survey, satshield}。

在计算维度，PISA 流水线的每个 MAU 级仅能在单个时钟周期内执行有限的 ALU 操作：哈希计算、整数加减法、位移、按位逻辑运算（AND/OR/XOR）。不支持浮点运算、乘除法、条件循环（\texttt{for/while}）或递归调用。这意味着任何需要迭代收敛（如梯度下降）或复杂数学运算（如对数、三角函数）的算法均无法在数据面直接实现~\cite{bosshart2013rmt}。

上述约束共同构成了本文算法设计的"硬边界"：所有在网检测逻辑必须是单遍（Single-pass）的、无循环的，且状态空间消耗必须严格控制在 MB 量级以内。这一约束直接驱动了第三章中“CMS 频率估计 + 固定容量候选缓存 + Bitmap”的级联两级检测架构设计~\cite{satshield}。


% =========================
% 扩展块 C：新增子节——关键链路形成机制与安全脆弱面
% 建议插入位置：第 40 行后（进入下一节前）
% =========================
\subsection{LEO 场景下关键链路形成机制与安全脆弱面}

巨型星座在几何意义上具有规则、稀疏且低度数的连接结构，但在业务承载意义上却会自然演化出一批“关键链路”。其形成原因可从三方面理解。首先，端到端业务需求矩阵本身具有显著的空间不均匀性：跨洲际流量、热点城市对业务、海洋与偏远区域回传业务以及若干大型信关站附近的聚合流量，会持续拉高特定区域路径的使用频度。其次，低时延优先的路由策略往往偏好更短跳数、更小传播距离或更少地面绕行的路径，这会使部分跨轨道或跨区域链路在大量业务对之间反复被选中。最后，极区异轨链路的周期性中断会阶段性削弱局部的路径多样性，迫使更多流量通过仍然可用的替代链路中继，进一步放大热点效应~\cite{handley2018delay_not_option, handley2019delay, laser_isl}。

为了描述这一长期趋势，可定义链路 $e$ 在连续 $K$ 个快照内的平均压力为
\begin{equation}
    \bar{\rho}_e=\frac{1}{K}\sum_{k=1}^{K}\rho_e(\tau_k),
\end{equation}
其中 $\rho_e(\tau_k)$ 为快照 $\tau_k$ 时刻链路利用率。若某条链路在多个快照上始终表现为高 $\bar{\rho}_e$、高路径承载比和低替代路径弹性，则该链路就构成了典型的结构性脆弱点。此类脆弱点的危险之处在于：它们并不一定长期拥塞，因此在日常运行中可能只表现为“稳定繁忙”；但一旦遭遇 LFA 这类刻意利用路径汇聚效应的攻击，其剩余带宽裕度会被快速吞噬，从而率先失稳。

这一机制与地面骨干网中的“核心边”概念相似，但 LEO 场景又具有两点额外强化因素。第一，星座拓扑的可预测性使得攻击者能够在快照尺度上提前推测热点链路的演化轨迹，而无需像地面互联网那样面对高度不透明的域间路由。第二，LEO 星座的平均节点度数较低，意味着单条链路失效或拥塞后，剩余替代路径的跳数膨胀与负载转移会更加明显，从而可能引发局部链路的级联拥塞。也就是说，LFA 在 LEO 中不仅能够“打满一条链路”，还可能通过触发绕行效应导致相邻链路的次级退化。

对本文而言，这一分析至少产生三个设计启示。其一，检测系统应重点关注那些位于高汇聚路径上的本地异常，而不是试图均匀对待所有链路；其二，协同系统的传播范围应围绕潜在瓶颈链路的上下游进行裁剪，以在有限带宽下获得最大的证据价值；其三，缓解系统不能只在受害点末端丢弃，而应尽可能将压制动作前推至热点链路之前的入口处，从源侧削减无效占用。这三点分别对应了第三章的候选监测思想、第四章的 ROI 设计和推回缓解思想。

% 建议新增图 2-3（待绘制后插入）
% 图题建议：LEO 星座中关键链路的形成与 LFA 攻击目标示意
% 画法建议：上半部分画正常业务在若干低时延路径上的自然汇聚，下半部分叠加 bot-decoy 攻击流，
%            以红色高亮 target link，并标出上游入口、受害链路与替代路径。
% \begin{figure}[t]
%   \centering
%   \includegraphics[width=0.95\linewidth]{figs/ch2/fig2_lfa_target_link.pdf}
%   \caption{LEO 星座中关键链路的形成与 LFA 攻击目标示意。}
%   \label{fig:ch2_lfa_target_link}
% \end{figure}

\section{链路洪泛攻击的攻击模型与对抗分析}

\subsection{DDoS 攻击的演进与 LFA 的定位}

分布式拒绝服务（Distributed Denial of Service, DDoS）攻击的演进可划分为三个世代。第一代为体积型攻击（Volumetric Attacks），以 UDP 反射放大（如 DNS/NTP/Memcached 反射）和 SYN Flood 为代表，其目标是以绝对的流量体量耗尽受害端的接入带宽或处理能力。此类攻击特征明显（流量突增、协议异常），可被基于简单阈值和五元组黑名单的流量清洗中心有效拦截。第二代为应用层攻击（Application Layer Attacks），如 HTTP GET/POST Flood 和 Slowloris，其流量体量较小但精准打击应用层资源（如 Web 服务器的并发连接数或数据库查询能力）。此类攻击需要更细粒度的应用层检测逻辑，但攻击流量仍然汇聚于受害终端，终端侧的异常可被直接感知~\cite{leosecurity}。

LFA 代表了第三代攻击范式的质变。其根本性的差异在于：攻击的受害者不再是任何终端节点，而是网络内部的中间链路（Target Link）。由于攻击流量的源地址和目的地址均为合法地址，且每条流的速率和行为模式与正常业务无异，终端侧（无论是 Bot 还是 Decoy）均不会感知到异常。唯一的"受害现场"是拥塞的中间链路，而这恰恰是传统终端侧安全设备的监控盲区。这一特点使 LFA 成为当前已知隐蔽性最高的网络层攻击手段之一~\cite{crossfire, coremelt}。

% [INS-ICARUS-05]
在 LEO 语境下，评估 LFA 攻防效果通常至少需要两个互补指标：一是攻击成本（达到目标拥塞所需的总注入带宽），二是可检测性代理指标（如攻击导致的上行最大增量 maxUp）。二者分别对应“攻击者资源门槛”与“运维侧可观测异常幅度”，有助于避免仅用单一成功率指标得出片面结论~\cite{icarus}。

\subsection{Crossfire 攻击的构造模型}

Crossfire~\cite{crossfire} 是 LFA 的经典实例，其攻击过程可形式化为三阶段模型。

\textbf{阶段一：拓扑测绘与瓶颈链路识别。} 攻击者控制分布在全球各地的僵尸主机集合 $\mathcal{B} = \{B_1, \ldots, B_N\}$，从每个 $B_i$ 向目标区域附近的公共服务器集合 $\mathcal{D} = \{D_1, \ldots, D_M\}$（如 CDN 节点、DNS 根服务器）执行 Traceroute，从而获取 $N \times M$ 条路径信息。通过对这些路径的重叠分析，攻击者可构建出一张目标网络的链路使用频率图（Link Usage Frequency Map），并从中识别出边介数中心性（Edge Betweenness Centrality）最高的链路子集 $E_{\text{target}}$。对于任意链路 $e$，其边介数中心性定义为：
\begin{equation}
    \text{BC}(e) = \sum_{(i,j)} \frac{\sigma_{ij}(e)}{\sigma_{ij}}
\end{equation}
其中 $\sigma_{ij}$ 为 $B_i$ 到 $D_j$ 的最短路径总数，$\sigma_{ij}(e)$ 为其中经过链路 $e$ 的路径数。$\text{BC}(e)$ 值越高，表明该链路承载的路径汇聚度越大，则为 LFA 目标~\cite{crossfire}。

\textbf{阶段二：Bot-Decoy 通信对分配。} 攻击者需要为每个 Bot $B_i$ 分配一组目标 Decoy $\{D_{j_1}, D_{j_2}, \ldots\}$，使得所有 Bot-Decoy 通信对在目标链路 $e^* \in E_{\text{target}}$ 上的汇聚流量总和超过其容量。形式化地，攻击者需要求解以下可行性问题：
\begin{equation}
    \sum_{\substack{(i,j): \\ e^* \in \text{path}(B_i, D_j)}} r_{ij} \geq C_{e^*}
\end{equation}
其中 $r_{ij}$ 为 $B_i$ 向 $D_j$ 发送的流量速率。为了保持隐蔽性，每条流的速率需满足 $r_{ij} \leq r_{\text{thresh}}$（典型值为数十 KB/s 至数百 KB/s），使其不触发终端侧或中间节点的单流速率告警~\cite{crossfire}。

\textbf{阶段三：协同流量投送。} 攻击者通过 C\&C（Command and Control）信道同步启动所有 Bot 的流量发送。在 LEO 场景中，由于星座拓扑的快照切换周期 $\Delta\tau$ 通常为 10$\sim$30\,s，攻击者可将攻击持续时间与快照周期对齐，在每个快照时段内保持稳定的攻击流配置，随拓扑切换同步调整 Bot-Decoy 分配方案~\cite{crossfire, icarus}。

% [INS-ICARUS-06]
进一步地，若攻击者掌握拓扑、路由策略及 Bot 地理位置等先验信息，其 C\&C 指令可按快照序列提前异步下发，无需在每个时隙实时重算并广播全量命令。该特性降低了攻击侧控制信令暴露，并使“实时拦截 C\&C”在 LEO 场景中更具挑战性~\cite{icarus}。

\subsection{LEO 场景下 LFA 的退化策略}

在 LEO 环境中，攻击者可利用星座拓扑的公开性和规律性设计针对性的退化策略（Degrading Strategies），以规避各类防御机制。本文重点关注以下三种退化策略，它们直接驱动了第三章多签名检测方案的特征维度选择。

\textbf{退化策略一：低速微流化（Low-rate Dilution）。} 攻击者通过扩大僵尸主机规模 $N$，将总攻击带宽分摊至更多的 Bot，使每条流的速率 $r_{ij}$ 降至极低值。设目标链路容量为 $C_{e^*}$，正常背景流量为 $\lambda_{\text{bg}}$，则攻击成功所需的 Bot 总流量为 $\Lambda_{\text{atk}} = C_{e^*} - \lambda_{\text{bg}}$。若 $N$ 足够大，则 $r_{ij} = \Lambda_{\text{atk}} / (N \cdot \bar{M}) \to 0$（$\bar{M}$ 为每个 Bot 的平均 Decoy 数），使得基于聚合速率的 Heavy-Hitter 检测完全失效~\cite{crossfire, ripple}。

\textbf{退化策略二：多靶机分散（Fan-out/Fan-in Relaxation）。} 攻击者增大每个 Bot 的目标 Decoy 数量和 Decoy 的地理分散度（即放大 Fan-out），同时增大指向同一 Decoy 的 Bot 数量（即放大 Fan-in），使攻击流量在源端和目的端均呈高度分散分布。这种分散性破坏了基于目的 IP 熵变化或目的前缀聚合的检测假设~\cite{crossfire}。

\textbf{退化策略三：脉冲式滚动攻击（On-Off Rolling）。} 攻击者将攻击分为交替的"开"（On）和"关"（Off）阶段，在"开"阶段以全力发送攻击流量，在"关"阶段暂停发送，使链路利用率回落至正常水平。通过控制 On-Off 的周期与占空比，攻击者可以使基于时间窗口平均值的检测算法难以捕捉瞬时突增。此外，攻击者可以在不同的快照时段内切换目标链路，形成时间和空间上的双重游击~\cite{icarus, dosat}。

% [INS-ICARUS-07]
需要注意的是，关于“动态性是否必然放大攻击收益”，现有结论并不单向。ICARUS 的初步结果显示：利用路径切换形成脉冲叠加在原理上可行，但可利用窗口在多数场景下并不充足；相比之下，若攻击脉冲与链路负载自然波动（load surge）时刻重合，则可能以更低代价触发拥塞。该问题仍依赖更细粒度的包级分析，属于当前研究边界而非既定定论~\cite{icarus, rasti2015temporal_lensing}。

上述三种退化策略的共同效果是：使 LFA 攻击流量在单一特征维度上与正常业务流量难以区分。这一分析直接驱动了第三章中引入多维度特征签名（速率、Fan-in/Fan-out、持续性）的设计决策——任何单一维度均可被退化策略规避，但多维度的联合异常仍然可以被检测到~\cite{satshield}。


% =========================
% 扩展块 D：新增子节——攻防中的观测不对称与设计启示
% 建议插入位置：第 87 行后（进入 P4 部分前）
% =========================
\subsection{LEO-LFA 攻防中的观测不对称与设计启示}

在前述攻击模型与退化策略基础上，还需要进一步指出：LEO-LFA 的真正难点并不仅仅在于攻击“隐蔽”，更在于攻防双方拥有不同类型的可观测信息与控制能力，形成了显著的观测不对称（Observation Asymmetry）。这一不对称性主要体现在以下四个方面。

\begin{enumerate}
    \item \textbf{拓扑知识不对称。} 攻击者可以利用公开的轨道根数、已知的路由偏好和多轮探测结果，在快照尺度上重建全局近似拓扑；而单颗卫星在默认情况下只能稳定观测经过本节点的局部切片，缺乏同等粒度的全局视角~\cite{crossfire, icarus}。
    \item \textbf{统计粒度不对称。} 攻击者在组织攻击时掌握 Bot、Decoy 与目标链路之间的全链路映射关系，可以按路径粒度进行流量编排；防御方若仅依赖本地聚合速率，往往只能获得被某条链路“截面化”之后的投影统计，难以直接重构完整的攻击编队。
    \item \textbf{执行位置不对称。} 攻击者可从源侧发起流量，天然在链路上游拥有动作发起点；而防御方若只在受害链路附近检测与丢弃，则往往只能在攻击流已经穿越多跳 ISL 之后被动响应，导致大量无效占用已经产生。
    \item \textbf{时间尺度不对称。} 攻击者可以依据快照切换和负载波动调整注入节奏，而防御方若依赖中心收集、离线分析或大窗口平均，便会在时间上持续滞后于攻击侧的重构速度。
\end{enumerate}

正是这种观测不对称，使得“单一特征 + 单点缓解”在 LEO 环境中很难奏效。单一特征的问题在于，它只能在有限的局部切片上寻找一类异常，而攻击者完全可以通过多 decoy 分散、低速微流化和脉冲滚动等手段将该异常压平；单点缓解的问题在于，它只能在本地收缩损害，却不能沿路径方向消除上游带宽浪费。因此，本文在后续章节中分别从两个方向削弱这种不对称：第三章通过多签名检测提升局部切片中的可分性，第四章通过局部协同与逐跳推回扩大观测与执行范围。换言之，第三章解决的是“看得更准”，第四章解决的是“管得更前”。

从论文结构安排来看，这一子节还承担着承上启下的作用：它将第二章中的背景性概念分析凝聚为面向方法设计的具体问题约束，有助于读者自然理解为什么你的后续方案必须同时包含多维签名、局部协同和推回缓解三个要素，而不是停留在传统的速率阈值检测或末端丢弃层面。

\section{在网计算与 P4 可编程数据面架构}

\subsection{从 OpenFlow 到 P4：数据面可编程能力的演进}

软件定义网络（SDN）通过解耦控制平面与数据平面，首次实现了对网络转发行为的集中式可编程控制。其代表性南向接口协议 OpenFlow 定义了一组固定的匹配字段集合（如 12-tuple 和后续扩展的 OXM 匹配字段），控制器通过下发流表规则指示数据面对匹配的数据包执行预定义的动作（转发、丢弃、修改字段等）。OpenFlow 的根本局限在于其"协议依赖性"：可匹配的协议头字段由标准预先定义并硬编码于交换芯片中，无法支持自定义协议头的解析与处理。当需要处理新的协议（如 INT 遥测头）或自定义封装格式时，必须等待标准更新和芯片厂商的硬件支持，这严重制约了数据面的创新速度~\cite{mckeown2008openflow}。

P4（Programming Protocol-independent Packet Processors）语言~\cite{bosshart2014p4}的出现解决了这一根本性问题。P4 的设计理念基于两个核心原则：协议无关性（Protocol Independence）——程序员可以自行定义数据包的解析逻辑和匹配字段，交换芯片不再绑定于任何特定协议；以及目标无关性（Target Independence）——同一份 P4 程序可以通过不同的编译器后端编译为不同硬件平台（如 ASIC 交换芯片、FPGA SmartNIC、软件交换机 BMv2）的可执行代码。P4 程序定义了完整的数据包处理流程：从比特流中解析出报文头（Parser），经过多级匹配-动作处理（Match-Action Pipeline），最终将修改后的报文头重新组装为比特流（Deparser）。

这种"自顶向下"的编程范式使得网络运营者和研究者能够根据需求灵活定义数据面的处理逻辑，无需受限于芯片硬编码的功能集。对于本文的研究而言，P4 的协议无关性意味着：我们可以在 LEO 卫星的交换芯片上自定义用于多签名检测的协议头部字段和状态更新逻辑，而无需依赖芯片厂商的专门支持~\cite{bosshart2014p4, p4_survey}。

\subsection{PISA 架构与匹配-动作流水线}

P4 程序运行于协议无关交换架构（Protocol Independent Switch Architecture, PISA）之上~\cite{bosshart2013rmt}。PISA 的硬件流水线由以下四个功能模块依次级联构成：

\textbf{（1）可编程解析器（Programmable Parser）。} 解析器是一个有限状态机（Finite State Machine, FSM），负责将输入的原始比特流按照 P4 程序定义的协议格式依次提取（Extract）各层报文头字段，并将提取结果写入包元数据（Packet Metadata）和包头向量（Parsed Header Vector, PHV）中。解析器支持条件分支（如根据 EtherType 字段决定下一层协议的解析方式），但不支持循环——每个数据包的解析路径必须是有向无环的。解析器的状态数和转换规则受硬件资源限制，典型 PISA 芯片支持的最大解析深度约为 $\sim$1000\,Bytes。

\textbf{（2）入口匹配-动作流水线（Ingress Match-Action Pipeline）。} 由 $S_{\text{ing}}$ 个（典型值 12$\sim$20）MAU 级（Stage）级联组成~\cite{bosshart2013rmt, p4_survey}。在每个 MAU 级 $s$（$1 \leq s \leq S_{\text{ing}}$），数据包的 PHV 中的指定字段被用作匹配键（Match Key），在该级的匹配表（Match Table，实现为 TCAM 或 SRAM 哈希表）中进行查找。匹配成功后，执行对应的动作（Action），包括修改 PHV 字段、读写有状态存储元素（下文详述）、设置出口端口等。每个 MAU 级的处理在单个时钟周期内完成，且各级之间严格串行，不存在回路或反馈——这是 PISA 实现线速处理的根本约束。

\textbf{（3）出口匹配-动作流水线（Egress Match-Action Pipeline）。} 结构与入口流水线类似，由 $S_{\text{egr}}$ 个 MAU 级组成，在数据包经过排队调度（Queuing and Scheduling）后执行。出口流水线通常用于数据包的最终修改（如更新校验和、添加遥测字段）。

\textbf{（4）逆解析器（Deparser）。} 将经过流水线处理后的 PHV 重新序列化为比特流，写入数据包的报文头区域，完成数据包的重组后发送至物理端口。

PISA 架构的核心约束可总结为：\textbf{无循环（No Loop）}——数据包在流水线中严格单向前进，不可回退；\textbf{固定时延（Deterministic Latency）}——每个数据包的处理时间恒定，不依赖于包内容或流状态；\textbf{有限级数（Bounded Stages）}——可用的 MAU 级数有限，复杂的处理逻辑需在有限级数内完成~\cite{bosshart2013rmt, p4_survey}。这些约束直接影响了第三章中多签名提取算法的流水线映射设计。

\subsection{数据面有状态处理与存储约束}

在 PISA 流水线中实现安全检测需要维护跨包的统计信息，这依赖于 P4 提供的有状态存储原语（Stateful Storage Primitives）。P4 定义了三类基本有状态元素~\cite{bosshart2014p4, p4_survey}：

\textbf{Register}：可在 MAU 的 ALU 中被读取和更新的整数数组。Register 是 P4 中最灵活的有状态存储单元，支持读-修改-写（Read-Modify-Write, RMW）操作，如递增、赋值、条件更新等。但存在重要约束：同一个 Register 数组的同一个索引（Index）在同一个流水线通过中只能被一个 MAU 级访问，以避免读写竞争（Read-Write Hazard）。

\textbf{Counter}：仅支持递增操作的计数器，适用于统计匹配特定规则的数据包数量或字节数。

\textbf{Meter}：实现令牌桶（Token Bucket）算法的速率测量器，可对流量进行三色（Green/Yellow/Red）标记，用于速率限制和流量整形。

上述有状态元素均驻留在 PISA 芯片的片上 SRAM 中。如前文所述，宇航级芯片的片上 SRAM 容量严格受限，通常为数 MB 量级。若第一级采用一个 $d\times w$ 的 CMS 计数矩阵（每个 counter 占 $b_{\text{ctr}}$ bit），并在第二级仅对 $K$ 个显式候选条目维护精细多签名状态（每个候选条目占用 $b$ 字节），则状态空间总消耗可近似写为
\begin{equation}
M_{\text{state}} \approx \frac{dwb_{\text{ctr}}}{8} + K \times b\ \text{Bytes}.
\end{equation}
以第三章的代表性部署口径为例，若取 $d=4,w=2048,b_{\text{ctr}}=16$ bit，则 CMS 矩阵约占 $16$~KB；若 $K=1024,b=128$，则候选级多签名状态约占 $128$~KB；再加上队列映射、双缓冲控制位与辅助寄存器，总预算仍处于数百 KB 量级，与第三章约 $200$~KB 的统一部署预算保持一致。这正是两级级联架构的设计优势：通过第一级的 sketch 频率估计与候选过滤，将需要进行精细多签名提取的流数量从全量（$\sim 10^6$）压缩至 $K$（$\sim 10^3$），从而将第二级的状态空间消耗控制在可接受范围内~\cite{satshield}。

此外，PISA 的 ALU 所支持的运算类型严格受限于整数加减法、位移、按位逻辑运算和哈希计算。不支持浮点运算、乘法（部分高端芯片支持有限的硬件乘法器）、除法和索引间接寻址（Indirect Indexing）。这一计算约束要求所有在网算法的数学操作必须可被分解为上述基本操作的组合。例如，第三章中的可疑度评分函数需要避免使用除法和指数运算，转而采用基于位移的近似计算或预计算查找表~\cite{bosshart2013rmt}。


% =========================
% 扩展块 E：新增子节——数据面、控制面与星上 agent 的职责划分
% 建议插入位置：第 129 行后（进入近似数据结构前）
% =========================
\subsection{数据面、控制面与星上 agent 的职责划分}

在介绍 PISA 流水线的能力边界之后，还需进一步回答一个工程实现层面的关键问题：哪些任务应当留在数据面逐包执行，哪些任务应当交由星上 agent 或地面控制系统完成？这一划分对于后续方案的可实现性具有决定性意义。一个常见误区是试图将所有检测、聚合与决策逻辑都压入 P4 程序中，但这既会迅速耗尽有限的流水线级数与寄存器资源，也会使程序难以适应快照切换、阈值重配置和协同消息处理等慢时间尺度任务。更合理的做法是采用“快路径在数据面、慢路径在 agent、长期优化在地面”的三层职责划分。

\begin{table}[t]
\centering
\caption{LEO 在网防御体系中不同执行层的职责划分}
\label{tab:ch2_plane_roles}
\small
\begin{tabular}{|p{2.4cm}|p{2.2cm}|p{4.0cm}|p{4.0cm}|}
\hline
执行层 & 时间尺度 & 典型任务 & 在本文中的对应功能 \\
\hline
数据面（P4/PISA） & ns--$\mu$s / 每包 & 报文解析、哈希、寄存器更新、位图置位、队列标记、TTL 递减 & 第三章中的 CMS 更新、候选缓存刷新、FI/FO 位图置位，以及第四章中的本地轻量消息处理 \\
\hline
星上 agent & ms--s / 每 epoch & 快照切换、阈值装载、位图读取、聚合评分、协同消息合并、推回软状态维护 & 第三章中的分数计算和 rank 映射更新，第四章中的 EG/PT 处理、抑制判定与推回安装 \\
\hline
地面控制 / 离线管理 & min--hour & 轨道/路由规划、FIB 预计算、长期参数调优、软件升级、策略审计 & 快照路由生成、长期统计分析、参数基线生成与离线评估 \\
\hline
\end{tabular}
\end{table}

上述三层划分的本质，是将不同时间尺度、不同复杂度和不同可信假设下的任务放置到最合适的执行位置。数据面负责所有必须“逐包、常数复杂度、立刻生效”的动作，例如计数更新、位图置位与队列标记；星上 agent 负责那些允许按窗口或按快照处理、但仍要求驻留卫星本地闭环的任务，例如分数计算、消息聚合与推回规则下发；而地面控制系统则只承担长期规划与离线优化，避免重新落回“每次异常都要走星地闭环”的旧模式~\cite{bosshart2013rmt, bosshart2014p4, p4_survey}。

这一职责划分同样解释了本文为何采用“数据面 + 星上 agent”协同架构，而非单纯强调某一侧的能力。若完全依赖数据面，则难以承载对数计算、窗口级读出与协同消息聚合；若完全依赖 agent 或地面控制，则又会牺牲包级响应的实时性。两者结合，才构成适合 LEO 场景的轻量级在网计算闭环。

% 建议新增图 2-4（待绘制后插入）
% 图题建议：数据面、星上 agent 与地面控制三层协同关系
% 画法建议：以一个卫星节点为中心，画 Parser/Ingress/Register/Egress 数据面快路径，
%            上方画星上 agent，左下画地面网络管理系统，箭头标注“逐包更新 / 每 epoch 读出 / 离线下发”。
% \begin{figure}[t]
%   \centering
%   \includegraphics[width=0.95\linewidth]{figs/ch2/fig2_plane_roles.pdf}
%   \caption{数据面、星上 agent 与地面控制三层协同关系。}
%   \label{fig:ch2_plane_roles}
% \end{figure}

\subsection{面向在网安全的近似数据结构}

上述存储与计算约束使得精确的 Per-flow 状态维护在星载环境中不可行，驱动了近似数据结构（Approximate/Probabilistic Data Structures）在在网安全检测中的广泛应用。本文涉及的核心近似数据结构包括以下四种。

\textbf{（1）布隆过滤器（Bloom Filter, BF）~\cite{bloom1970space}。} BF 由一个长度为 $m$ 比特的位数组（Bit Array）和 $k$ 个独立的哈希函数 $h_1, h_2, \ldots, h_k$ 组成。插入元素 $x$ 时，将位数组中 $h_1(x), h_2(x), \ldots, h_k(x)$ 对应的位置置为 1。查询元素 $y$ 时，检查 $h_1(y), \ldots, h_k(y)$ 对应的位是否全为 1。BF 不存在假阴性（False Negative），但存在假阳性（False Positive），假阳性率约为 $(1-e^{-kn/m})^k$，其中 $n$ 为已插入元素数量。本文第四章的协同消息去重与 seen-set 维护可采用 BF 实现，从而在局部协同范围内以常数级状态开销实现快速查重。

\textbf{（2）Count-Min Sketch（CMS）~\cite{cormode2005count}。} CMS 由 $d$ 行 $w$ 列的二维计数器矩阵组成，配合 $d$ 个独立的哈希函数 $h_1, \ldots, h_d$。更新元素 $x$ 的计数时，对每行 $i$，将 $\text{CMS}[i][h_i(x)]$ 递增 1。查询 $x$ 的频次时，返回 $\hat{f}(x) = \min_{1 \leq i \leq d} \text{CMS}[i][h_i(x)]$。CMS 的估计具有单边误差性质：$\hat{f}(x) \geq f(x)$（始终高估），且 $\Pr[\hat{f}(x) - f(x) > \epsilon N] < \delta$，其中 $N$ 为总元素数，$\epsilon = e/w, \delta = e^{-d}$。CMS 的优点在于：每个包仅需执行固定次数的哈希与计数更新，天然适配 P4 流水线的常数复杂度约束。其局限在于只保存频次矩阵而不显式保存 key 身份，因此若要恢复 Top-$k$ 候选，还需要与显式候选缓存或堆结构配合。本文第三章即采用“CMS + 固定容量候选缓存”的组合作为一级候选生成器。

\textbf{（3）固定容量候选缓存（Candidate Cache）。} 候选缓存由容量为 $K$ 的显式条目表组成，每个条目保存候选 key 的标识、当前估计频次以及与后续二级特征提取相关的元数据（如位图槽位索引、计数寄存器索引等）。与纯 CMS 不同，候选缓存显式保留了 key 身份，因此可为后续的 FI/FO 位图、Persist 统计和队列映射提供可索引对象；与需要维持全局严格有序列表的 exact Top-$k$ 排序不同，候选缓存只需在固定容量下执行“命中则刷新、未命中则按阈值晋升/替换”的局部更新，更便于在 P4 数据面或 agent 协同控制中实现。第三章正是通过“CMS 负责全局频次估计，候选缓存负责显式 key 恢复”的分工，近似恢复 Top-$k$ 候选集合。

\textbf{（4）位图估计（Bitmap / Linear Counting）~\cite{whang1990linear}。} 使用一个 $m$ 比特的位数组，通过哈希函数将元素映射到对应位并置位。集合基数的估计基于空位数量 $V$：$\hat{n} = -m \ln(V/m)$。位图估计在精度可控的前提下将基数统计问题转化为简单的位操作，天然适配 P4 的 ALU 约束。在第三章中，位图被用于 Fan-in 和 Fan-out 的在线近似统计——对候选缓存中的 key 维护独立的源 IP 位图和目的 IP 位图，每个数据包到达时仅需执行一次哈希计算和一次位置位操作。


% =========================
% 扩展块 F：近似数据结构的对比总结表与设计映射
% 建议插入位置：第 141 行后
% =========================

从工程映射角度看，上述四类近似数据结构并不是相互替代关系，而是共同构成了“粗筛—显式恢复—细粒度统计—消息去重”的功能链条：CMS 负责在全量流空间中给出低成本的频率估计，候选缓存负责恢复 key 身份并承接后续精细特征，Bitmap 负责在候选范围内提取 distinct 类结构特征，而 Bloom Filter 则适合承担第四章中 seen-set 这样的存在性判断任务。它们之所以适合在星载场景联合使用，恰恰在于每个包经过时只需要执行常数次哈希、读写和位操作，而无需任何复杂迭代或全表扫描。

\begin{table}[t]
\centering
\caption{本文所用近似数据结构的功能定位与实现代价对比}
\label{tab:ch2_ds_compare}
\small
\begin{tabular}{|p{2.6cm}|p{2.5cm}|p{2.3cm}|p{2.6cm}|p{3.5cm}|}
\hline
结构 & 核心操作 & 误差特性 & 状态/更新代价 & 在本文中的定位 \\
\hline
Bloom Filter & 插入、存在性查询 & 无假阴性，存在假阳性 & 位数组 + 常数次哈希 & 第四章协同消息去重、seen-set 维护 \\
\hline
Count-Min Sketch & 频次更新、点查询 & 单边高估 & $d$ 次计数器更新 & 第三章全量流空间中的一级频率估计 \\
\hline
Candidate Cache & 命中、刷新、替换 & 近似恢复 Top-$K$ 候选 & 固定容量显式状态 & 恢复 key 身份并承接二级特征统计 \\
\hline
Bitmap / Linear Counting & 置位、空位计数、基数估计 & 接近饱和时误差增大 & 一次置位 + 窗口末估计 & 第三章 FI/FO 的 distinct 近似统计 \\
\hline
\end{tabular}
\end{table}

这一表格建议保留在第二章中，因为它能够显著提升读者对“为什么第三章不是只用一种 sketch 就解决所有问题”的理解。更具体地说，LEO 在网安全防御并不缺少某一种单独的数据结构，而是缺少一条符合 PISA 约束的功能链：哪一层负责全量监视，哪一层负责显式 key 恢复，哪一层负责复杂结构特征，哪一层负责分布式协同中的消息控制。把这些功能链在第二章中讲清楚，第三章和第四章的具体实现就会显得顺理成章，而不是“突然出现一堆结构拼装”。

\section{分布式 Gossip 协议与协同防御基础}

\subsection{去中心化协同的必要性与传统共识的代价}

第三章的单星检测方案解决了本地闭环响应的问题，但 LFA 攻击的碎片化特性使得任何单颗卫星只能观测到途经本节点的局部流量切片。在退化策略下，每颗卫星本地观测到的单流速率可能完全正常，但多颗卫星上的观测结果汇聚后才能揭示出异常的汇聚模式~\cite{crossfire, satshield}。因此，多星之间的检测证据共享与协同判定是有效对抗 LFA 的必要条件。

直觉上，多节点协同可以通过集中式方案实现——即设置一个或多个中心协调节点，各卫星将本地检测结果上报至中心节点，由后者做出全局判定。但在 LEO 环境中，这一方案面临三个固有问题：（1）中心节点的指定与维护在高动态拓扑下本身就是一个非平凡问题（需要类似 Leader Election 的机制）；（2）中心节点到各卫星的通信路径可能跨越多个 ISL 跳数，引入额外时延；（3）中心节点本身可能是 LFA 攻击的目标，或其接入链路可能被攻击者瘫痪~\cite{crossfire}。

另一种方案是采用分布式强一致性共识协议（如 Paxos、Raft），使所有参与节点就攻击判定达成一致。但此类协议的通信复杂度和时延代价在 LEO 环境中不可接受：以 Raft 为例，一次共识需要至少两轮消息交换（提案-确认），涉及过半节点的参与。在 $n$ 个节点组成的集群中，单次共识的消息复杂度为 $O(n)$，如果区域内有数十颗卫星参与协同，每次判定需要交换数十至数百条控制消息。在 ISL 带宽已被 LFA 攻击流量大量侵占的情况下，这些控制消息本身可能因拥塞而丢失或严重延迟，导致协议无法收敛。

全网广播泛洪（Flooding）是最简单的信息共享方式，但其消息复杂度为 $O(|V| \cdot |E|)$，在大规模星座（$|V| > 1000$）中将产生灾难性的信令风暴（Broadcast Storm），不仅耗尽剩余的可用带宽，还可能触发二次拥塞。

\subsection{Gossip 协议的基本原理与传播特性}

Gossip 协议（又称流行病协议, Epidemic Protocol）~\cite{demers1987epidemic}提供了一种介于强一致性共识与盲目泛洪之间的折中方案。其基本机制为：每个节点周期性地（周期为 $T_{\text{gossip}}$）随机选取一个或少数几个邻居节点，将本地持有的信息（如攻击检测证据）发送给对方；接收方将收到的信息与本地信息进行合并（Merge），并在后续周期中继续向其邻居传播。

Gossip 协议的传播动力学可类比为传染病模型中的 SIR（Susceptible-Infected-Recovered）模型。设网络中有 $n$ 个节点，初始时有 $i_0$ 个节点持有某条信息（"已感染"），每个已感染节点在每个 Gossip 周期中以概率 $p$ 将信息传播给一个随机邻居。在完全图（Complete Graph）上，信息到达所有节点所需的轮次为 $O(\log n)$，消息总量为 $O(n \log n)$。在稀疏图（如 LEO 的 4-regular mesh）上，传播速度取决于图的连通性和度数，但仍然保持对数级别的轮次复杂度~\cite{demers1987epidemic}。

Gossip 协议的关键优势在于：（1）\textbf{去中心化}——无需指定或维护中心协调节点；（2）\textbf{容错性}——即使部分节点或链路失效，信息仍可通过替代路径传播，最终到达所有可达节点（最终一致性, Eventual Consistency）；（3）\textbf{可扩展性}——每个节点在每个周期中仅与少量邻居通信，单节点的通信开销与网络规模无关。这些特性使 Gossip 协议天然适合在高动态、部分链路可能被攻击瘫痪的 LEO 环境中进行多星协同~\cite{demers1987epidemic}。

\subsection{抑制型传播与开销控制}

标准 Gossip 协议的一个固有问题是冗余消息的过度生成：已经获知某条信息的节点仍然可能反复收到相同信息的推送，造成带宽浪费。在 LFA 攻击已导致 ISL 带宽紧张的环境中，这种冗余尤其不可接受。

抑制型传播（Suppressed Broadcasting）机制通过引入冗余检测与主动静默来控制 Gossip 的开销。其核心思想源自物联网领域的 Trickle 算法~\cite{levis2004trickle}，可形式化为以下流程。

对于每条待传播的消息 $m$，节点维护三个状态变量：传播计时器 $T_m$（在区间 $[T_{\min}, T_{\max}]$ 内动态调整）、冗余计数器 $c_m$（记录在当前周期内收到的与 $m$ 内容一致的消息数量）、以及冗余抑制阈值 $K$。在每个传播周期开始时，节点将 $c_m$ 重置为 0 并启动计时器 $T_m$。在计时器运行期间，每当节点从邻居收到一条与 $m$ 内容一致的消息时，将 $c_m$ 递增 1。当计时器到期时，若 $c_m < K$，节点判断其邻域内尚未充分获知该信息，遂主动向邻居发送 $m$；若 $c_m \geq K$，节点判断其邻域已充分获知该信息，遂抑制本次发送。若信息在一段时间内未再听到新的更新，$T_m$ 倍增至 $T_{\max}$，节点进入静默状态~\cite{levis2004trickle}。

此外，通过为每条消息设置 TTL（Time-To-Live）跳数限制，可将其传播范围限制在发起节点的 $h$-hop 邻域内（即感兴趣区域 Region of Interest, ROI），避免信息向无关区域扩散。在第四章的 CoopSatShield 设计中，ROI 的跳数半径 $h$ 被设定为与 LFA 攻击路径的典型跳数相匹配，确保协同判定的参与节点恰好覆盖攻击路径经过的卫星群组~\cite{satshield}。


% =========================
% 扩展块 G：新增子节——Gossip 参数选择与区域协同范围设计
% 建议插入位置：第 171 行后（本章小结前）
% =========================
\subsection{Gossip 参数选择与区域协同范围设计}

在将 Gossip 协议用于 LEO 多星协同时，真正决定系统表现的并不仅是“是否采用 Gossip”，而是具体的参数配置。与地面大规模互联网环境不同，LEO 场景中的协同传播必须同时服从三个目标：其一，在攻击持续期间尽快把局部证据扩散到足够多的相关节点；其二，将控制消息限制在真正与受害路径有关的局部区域内；其三，在带宽已经紧张的 ISL 上避免因冗余传播而制造新的拥塞。因此，TTL 半径 $h$、最小/最大计时器 $T_{\min},T_{\max}$ 以及冗余阈值 $K$ 不是普通的实现细节，而是决定协同系统“快、准、省”三者平衡的核心控制旋钮~\cite{demers1987epidemic, levis2004trickle}。

\begin{table}[t]
\centering
\caption{LEO 多星协同候选机制的对比}
\label{tab:ch2_collaboration_compare}
\small
\begin{tabular}{|p{2.2cm}|p{2.3cm}|p{2.1cm}|p{2.1cm}|p{3.2cm}|}
\hline
机制 & 组织方式 & 典型时延 & 信令开销 & 对 LEO 的适配性 \\
\hline
中心化汇聚 & 单中心或少量中心节点收集证据 & 中到高 & 中等 & 易受长路径时延、单点故障与中心链路拥塞影响 \\
\hline
强一致性共识 & 多节点多数派确认后再决策 & 高 & 高 & 一致性强，但在高动态和拥塞环境下代价偏大 \\
\hline
全网广播泛洪 & 所有节点无差别扩散消息 & 低到中 & 极高 & 实现简单，但极易引发广播风暴 \\
\hline
抑制型 Gossip & 局部随机扩散 + 冗余抑制 + TTL 裁剪 & 中低 & 低到中 & 去中心化、容错性好，适合 LEO 局部协同 \\
\hline
\end{tabular}
\end{table}

一般而言，$h$ 主要决定协同的空间边界：$h$ 过小，证据传播无法覆盖攻击路径上的关键上游节点；$h$ 过大，则会让无关卫星被卷入协同，平白增加控制开销。$T_{\min}$ 决定首轮扩散速度，应足够小以追上攻击状态变化，但又不能小到使多个节点同时争抢发送、造成过多碰撞；$T_{\max}$ 决定稳定期的静默程度，通常用于在没有新证据时逐步降低控制频率。冗余阈值 $K$ 则体现“本地已经足够知道这件事了吗”的判断：较小的 $K$ 会更积极地抑制重复消息，较大的 $K$ 则更偏向于保证覆盖度。若以 $h$ 跳局部区域为协同边界，且每轮扩散最短受 $T_{\min}$ 限制，则消息覆盖到 $h$ 跳邻域的时间通常处于 $O(hT_{\min})$ 的量级；这一定性关系也解释了为什么第四章中需要在覆盖范围与传播时延之间寻找平衡。

对本文方案而言，一个重要的设计原则是：Gossip 不应被当作“全网通告协议”，而应被当作“受害区域附近的受控证据聚合机制”。也就是说，它的目标不是让所有卫星都知道某条证据，而是让那些最可能共享攻击路径、最有机会执行推回动作的卫星尽快知道这条证据。正是在这一原则下，第四章采用了“局部区域限制 + 冗余抑制 + 软状态寿命控制”的组合设计，从而将协同开销控制在远低于业务带宽的范围内，同时又保证了证据能够沿最有价值的局部方向扩散。

% 建议新增图 2-5（待绘制后插入）
% 图题建议：抑制型 Gossip 在 ROI 内的传播与静默机制
% 画法建议：画一个 h-hop ROI，起点为受害链路附近节点，消息以多条边扩散；
%            对已听到重复证据的节点打“mute”标记，并用虚线表示被抑制的发送。
% \begin{figure}[t]
%   \centering
%   \includegraphics[width=0.9\linewidth]{figs/ch2/fig2_gossip_roi.pdf}
%   \caption{抑制型 Gossip 在 ROI 内的传播与静默机制。}
%   \label{fig:ch2_gossip_roi}
% \end{figure}

\section{本章小结}

本章从四个维度构建了支撑后续研究的理论基础。在网络架构层面，建立了 LEO 巨型星座的时变图模型 $G(t)=(V(t), E(t))$ 和快照路由模型，明确了 ISL 容量约束 $C_e$ 和拓扑切换周期 $\Delta\tau$ 的量化参数。在威胁模型层面，形式化描述了 Crossfire 型 LFA 的三阶段攻击流程和三种退化策略，论证了多维度特征检测的必要性。在计算平台层面，详细分析了 PISA 流水线的无循环、固定级数、有限 ALU 等硬约束，以及片上 SRAM 的容量上限，建立了在网算法设计的"硬边界"条件。在协同机制层面，引入了 Gossip 协议的传播理论与抑制型变种的开销控制模型，论证了其在 LEO 高动态、部分链路受损环境下的适用性。这些理论基础将在第三章（多签名在网检测）和第四章（多星协同推回）中被直接引用和具体化。
